\section{Discussion and Limitations}
\label{sec:discussion}

The two simulators expose different parts of the inverse problem. Eight voltage/current PMUs recover the requested 31-bus voltage field accurately in the PowerDynamics tests even though the 114-state local model is not fully observable. The same sensor buses support reliable abnormal-event detection and moderate known-source localization in ANDES. Neither result transfers to the other task. Low hidden-voltage HVRE does not guarantee source separation, and high detection F1 does not identify a hidden state or a physical source.

The local theory explains this split. Proposition~\ref{prop:functional-reconstruction} removes unobservable directions that do not affect the requested voltage function. Corollary~\ref{cor:functional-blue} then attaches a covariance to that estimable function; rank alone cannot provide it. Source diagnosis has a different geometry. Proposition~\ref{prop:source-indistinguishability} tests whether two candidate-nuisance response sets intersect, while Proposition~\ref{prop:nonlinear-separation} requires the affine distance to exceed the combined DAE remainder before local nonlinear separation is certified. Large faults, topology switches, and unmodeled controls may violate the neighborhood in which that bound holds.

Only the first theory-to-data link has been tested. The preregistered functional residual correlates with bus-level error, although excitation, scaling, estimator covariance, and electrical distance explain several exceptions. The source-held-out failure is compatible with poor physical separation, source memorization, or both. The current artifacts cannot separate those explanations. Computing \eqref{eq:pairwise-information} before model fitting and relating it to held-source errors would test the mechanism directly.

The operating-point study also separates target recovery from parameter recovery. Regularized physical recentering closes \EZeroHClosureOverall{} of the oracle-center gap on M6 although only \EZeroHRank{} of \EZeroHDimension{} nuisance directions are locally informed. The hidden-voltage target can therefore be accurate while the injection coordinates remain non-unique. This result is confined to an equilibrium snapshot. It does not cover line outages, faults, or sequential event-conditioned estimation, and the local curvature covariance remains under-dispersed in the joint NEES-like statistic.

The legacy event model supplies one reusable mechanism and one clear warning. Multiscale summaries, adaptive residuals, and simultaneous PMU contrasts improve known-source physical ranking by \LegacyPhysicalGainPP{} percentage points. Availability belongs in the integrity state because a missing measurement is observed directly. Source identities and endpoint descriptors, however, allow memorization. Exact localization falls to \HeldSourceRate{} under complete physical-source holdout. A source-generalizable model should instead compute engineered features from candidate residuals, $\bm y-\widehat{\bm y}_h$, and compare their contribution with a physics-only likelihood.

The principal evidence gap is the absence of one end-to-end experiment. The state and event studies use different simulators, measurement semantics, and split rules. Both are simulated; neither supplies field or hardware-in-the-loop validation. The event records also have fixed onsets and only about 7.15~min of normal-labeled exposure per split, too little for an operational false-alarm rate. PMU loss and addition have not been linked prospectively to the separation margin, and the candidate posterior has not been calibrated.

A decisive follow-up requires one audited plant and measurement operator for learned-only, physics-only, learned-residual, and diagnosability-aware methods. Event parents and physical sources must be disjoint across fitting, calibration, and test. Operating-point shift, parameter error, and PMU loss or addition should be imposed before the test is opened. A long normal stream must report false alarms per hour and event recall; event-source correctness, candidate-set coverage, resolution delay, HVRE, calibration error, runtime, and failed solves complete the evaluation. Without that experiment, joint superiority remains unproved.
